\chapter{Enteral Feeding Tube Placement}

\section*{What Is This Test or Procedure?}
Enteral feeding tube placement is the insertion of a tube into the gastrointestinal tract to deliver nutrition, fluids, and medications directly to the stomach or small intestine. Access may be short-term (nasogastric or nasoenteric tubes) or long-term (gastrostomy or jejunostomy).

\section*{Alternative Names}
Nasogastric (NG) tube placement, Nasojejunal (NJ) tube placement, Gastrostomy tube (G-tube), Percutaneous endoscopic gastrostomy (PEG), Jejunostomy tube

\section*{Principle}
Enteral feeding delivers formula directly into the gastrointestinal tract, preserving gut mucosal integrity and using the body's normal digestive and absorptive pathways. Nasoenteric tubes are passed through the nose into the stomach or small bowel, with position confirmed by pH of aspirate or radiography. Long-term access is created surgically or endoscopically by placing a catheter directly through the abdominal wall into the stomach or jejunum.

\section*{History}
Nasogastric feeding was described in the 19th century. Percutaneous endoscopic gastrostomy, introduced in 1980, became the standard approach for long-term enteral access.

\section*{Why This Test or Procedure Is Ordered}
Enteral feeding is ordered when a patient cannot eat enough by mouth but has a functioning gastrointestinal tract, as in dysphagia, stroke, major trauma, burns, severe malnutrition, or prolonged critical illness. It provides nutrition, hydration, and medications and may also be used for gastric decompression.

\section*{Is This a Primary or Secondary Test or Procedure}
Primary therapeutic procedure for delivering enteral nutrition when oral intake is inadequate or unsafe.

\section*{How This Test or Procedure Is Performed}
For nasoenteric tubes, the tube is lubricated and passed through the nose into the stomach, with advancement into the small bowel assisted by technique or image guidance. For gastrostomy, an endoscope identifies a safe site and the tube is placed percutaneously through the abdominal wall. Position is confirmed by aspirate pH or imaging before feeding begins.

\section*{What Preparation Is Needed}
Informed consent is obtained; nasoenteric placement usually requires no fasting, whereas gastrostomy requires fasting and often prophylactic antibiotics. Coagulation status is checked before percutaneous placement, and cervical spine clearance is needed before nasal passage in trauma patients.

\section*{Contraindications and Precautions}
Enteral feeding is avoided in bowel obstruction, severe ileus, or peritonitis. Nasoenteric tubes are avoided with facial trauma, skull base fractures, or esophageal strictures, and percutaneous placement is avoided with coagulopathy, ascites, or anatomic barriers. Use caution in patients at high risk for aspiration.

\section*{Risks}
Misplacement into the airway, aspiration pneumonia, epistaxis, and tube dislodgment or clogging. Percutaneous placement carries risks of bleeding, bowel perforation, peritonitis, and site infection.

\section*{What Happens After the Test or Procedure}
Tube position is verified before each use, and feedings are started slowly and advanced as tolerated. The site is cleaned daily, and the tube is flushed after each use to prevent clogging and infection.

\section*{Understanding Your Results}
Outcome is determined by tolerance of the feeding regimen and the absence of complications rather than by a numeric value. Success is judged by weight stabilization or gain and achievement of metabolic targets.

\section*{Normal Results}
There is no quantitative normal value; the goal is a well-tolerated feeding schedule, stable weight, and no evidence of aspiration or tube-related complications.

\section*{Follow-up Tests and Procedures}
Nutritional assessment and weight monitoring continue, with imaging if tube displacement is suspected. Long-term plans include transitioning back to oral feeding when safe or exchanging to a low-profile device.

\section*{References}
\begin{enumerate}
  \item MedlinePlus. Enteral feeding tube placement. U.S. National Library of Medicine; 2025.
  \item Current Medical Diagnosis and Treatment. McGraw-Hill; 2025.
\end{enumerate}

\clearpage
